\section{Related Work}

\paragraph{Personality and situated behavior.}
TRAIT extends personality inventories into scenario-based choices and evaluates
distinctiveness, consistency, and prompting effects \citep{lee2025trait}.
Behavioral Mode Axes uses situated choices across interaction registers,
also evaluates open-ended generations, and studies activation-based control
\citep{liu2026bma}. Separate construction and evaluation scenarios are already
part of that framework. Internal-feature interventions further connect
personality-related representations to behavior across situations
\citep{zhang2026triad}; conversational studies examine contextual variation under
the same assigned personality \citep{han2026contexts}. Together these studies
motivate examining both recurring defaults and conditional expression. Our
deployment-level forecasts test observable educational actions and do not
identify the internal mechanisms responsible for them.

\paragraph{Tutor behavior and educational personas.}
The CIMA comparison studies three LLMs in language-learning dialogues, reporting
characteristic action mixtures and response complexity. Its generation prompt
requests both a response and action labels \citep{kucheria2025cima}. Thus,
describing model-specific teaching styles is an existing research direction.
Tutor-persona steering learns from human tutoring dialogues and evaluates
persona alignment on held-out dialogues \citep{lee2026tutorpersonas}. Our
experiment instead measures the default behavior of multiple deployed
configurations and compares changes caused by controlled student cues and
teaching instructions. We code the visible answers after generation and retain
cross-coder sensitivity; this changes the measurement procedure without
establishing that the resulting labels are ground truth.

\paragraph{Teaching quality and the predictive question.}
MRBench and MathTutorBench organize pedagogical-capability evaluation
\citep{maurya2025mrbench,macina2025mathtutorbench}. PATS studies teaching adapted
to student personality \citep{rooein2026pats}. Our primary target is the
probability of the tutor's next observable behavior under specified conditions.
Predicting that behavior does not rank its educational value. The empirical
addition sought here is a joint account of which default and conditional
profiles predict new-source actions, how those profiles compare with simpler
explanations, and how neutral wording and explicit policies change them on
matched inputs. Source separation and advance prediction locks make these
comparisons auditable; their substantive contribution depends on the resulting
evidence rather than on the procedure alone.
